\chapter{Liaisons hertziennes longue distance}\label{ch:nw:long-haul-wireless}

Quand une tempête traverse la ligne de pylônes entre Chicago et le New Jersey, les \omterm{def:nw:long-haul-wireless:mw}{liaisons hertziennes micro-ondes} situées sous elle s'évanouissent, et tant que dure la pluie
la route la plus rapide entre les contrats à terme et les actions redevient la fibre, plus lente de trois millisecondes. Shkilko et Sokolov ont précisément
mis à profit ces après-midi comme expérience~: des épisodes météorologiques qui «~suppriment temporairement les avantages de vitesse des traders les plus rapides en perturbant
leurs réseaux micro-ondes~» étaient associés à une sélection adverse et à des coûts de transaction plus faibles. La pluie a donné aux traders les plus lents quelques
heures sans cette course, et la liquidité du marché s'est améliorée tant qu'elle a duré.

Le chapitre~10 a mesuré les routes à l'aune de leurs planchers, et le chapitre~13 a disséqué la fibre optique. Le présent chapitre fait de même pour l'air~: la géométrie
d'un saut en visibilité directe, la puissance dont il dispose en réserve, la pluie qui la lui retire, les \omterm{def:nw:long-haul-wireless:mmw}{liaisons en ondes millimétriques} et les liaisons laser au sein d'une
zone métropolitaine, la radio en ondes courtes qui franchit les océans, et ce que ces liaisons peuvent transporter. Son modèle fait traverser une tempête étiquetée une chaîne de pylônes
entre Aurora et Carteret et compte les minutes passées sur la fibre.

\section{Liaisons hertziennes micro-ondes~: trajet en visibilité directe, zones de Fresnel et pylônes}

\begin{definition}[Liaison hertzienne micro-ondes, trajet en visibilité directe]\label{def:nw:long-haul-wireless:mw}
Une \emph{liaison hertzienne micro-ondes}\index{liaison hertzienne micro-ondes} (microwave link) transporte des données par radio à des fréquences de quelques à quelques dizaines de gigahertz entre deux
antennes directionnelles, typiquement sur des pylônes ou des toits. Elle nécessite un \emph{trajet en visibilité directe}\index{trajet en visibilité directe} (line-of-sight path)~: un chemin rectiligne entre les antennes
dégagé du sol et des obstacles, car de telles ondes ne les contournent pas.
\end{definition}

\begin{definition}[Zone de Fresnel, renflement terrestre]\label{def:nw:long-haul-wireless:fresnel}
La première \emph{zone de Fresnel}\index{zone de Fresnel} (Fresnel zone) d'un trajet radio est l'ellipsoïde autour de la droite joignant les antennes, dans laquelle
les trajets sont au plus d'une demi-longueur d'onde plus longs que le trajet direct~; les obstacles situés à l'intérieur affaiblissent le signal reçu, si bien qu'une liaison est conçue pour la garder
dégagée. Le \emph{renflement terrestre}\index{renflement terrestre} (earth bulge) est la hauteur de la surface terrestre au-dessus de la droite joignant deux points au niveau du
sol, en un point donné entre eux~; la réfraction atmosphérique infléchit légèrement vers le bas les ondes radio, ce qui se modélise par une Terre de rayon
$K$ fois le rayon réel.
\end{definition}

\begin{proposition}[Géométrie du saut]\label{prop:nw:long-haul-wireless:geometry}
À des distances $d_1$ et $d_2$ kilomètres des deux extrémités d'un saut de longueur $d = d_1 + d_2$, le \omterm{def:nw:long-haul-wireless:fresnel}{renflement terrestre} vaut
$h = d_1 d_2 / (12.74\,K)$ mètres et le rayon de la première \omterm{def:nw:long-haul-wireless:fresnel}{zone de Fresnel} à $f$ gigahertz est $r = 17.32\sqrt{d_1 d_2 / (f\,d)}$ mètres. Au
milieu du saut, $h = d^2/(50.96\,K)$ et $r = 8.66\sqrt{d/f}$~: le renflement croît avec le carré du saut, la zone avec sa racine carrée.
\end{proposition}

\begin{proof}
Un cercle de rayon $K R_\oplus$ s'élève au-dessus de sa corde, à des distances $d_1$ et $d_2$ de ses extrémités, de $d_1 d_2/(2 K R_\oplus)$ au premier ordre~;
avec $R_\oplus = 6\,371$~km et le résultat en mètres, $2 R_\oplus/1\,000 = 12.74$. Le premier rayon de Fresnel vaut $\sqrt{\lambda d_1 d_2/d}$ avec
$\lambda = c_0/f$~; avec $d$ en kilomètres et $f$ en gigahertz, la constante vaut $\sqrt{0.2998 \times 10^3} = 17.32$. En posant $d_1 = d_2 = d/2$
on obtient les expressions au milieu du saut.
\end{proof}

Laughlin, Aguirre et Grundfest donnent les mêmes règles au milieu du saut (8.7 et 1/50) et en tirent la conséquence pour les réseaux Chicago--New York~: avec
des pylônes d'environ 100 mètres, des fréquences proches de \qty{6}{\giga\hertz}, $K$ inférieur à 4/3 et des obstacles de 10 mètres, un saut au-dessus d'un terrain plat ne peut guère
dépasser environ 70 kilomètres, si bien qu'une route de 1\,200 kilomètres exige vingt sauts ou plus. Le \Cref{tab:nw:long-haul-wireless:geometry}
reproduit le calcul.

\begin{table}[ht]
\centering\small
\begin{tabular}{rrrrr}
\toprule
Saut & Fréquence & \omterm{def:nw:long-haul-wireless:fresnel}{Renflement terrestre} & Rayon de Fresnel & Hauteur des pylônes \\
(km) & (GHz) & au milieu du saut (m) & au milieu du saut (m) & (m) \\
\midrule
30 & 6 & 13.2 & 19.4 & 42.6 \\
30 & 11 & 13.2 & 14.3 & 37.5 \\
50 & 6 & 36.8 & 25.0 & 71.8 \\
50 & 11 & 36.8 & 18.5 & 65.3 \\
70 & 6 & 72.1 & 29.6 & 111.7 \\
70 & 11 & 72.1 & 21.8 & 104.0 \\
\bottomrule
\end{tabular}
\caption{Géométrie du saut d'après la \cref{prop:nw:long-haul-wireless:geometry} avec $K = 4/3$~: la hauteur que doivent avoir deux pylônes identiques pour qu'un obstacle de 10~mètres
au milieu du saut reste dégagé de toute la première \omterm{def:nw:long-haul-wireless:fresnel}{zone de Fresnel}. Au-delà de 70 kilomètres, les pylônes dépassent 100 mètres. Données~:
\texttt{nw\_wireless.geometry\_rows()}.}\label{tab:nw:long-haul-wireless:geometry}
\end{table}

\begin{omfigure}
\begin{tikzpicture}[font=\footnotesize,>=Stealth]
\draw[thick,gray] (0,0) arc[start angle=110,end angle=70,radius=16.4];
\draw[thick] (0,0) -- (0,2.3); \draw[thick] (11.2,0) -- (11.2,2.3);
\node[below] at (0,0) {pylône A}; \node[below] at (11.2,0) {pylône B};
\draw[omDef,thick] (0,2.3) -- (11.2,2.3);
\draw[omThm,dashed] (5.6,2.3) ellipse (5.6 and 0.55);
\draw[gray,densely dotted] (0,0) -- (11.2,0);
\draw[<->,omMeth] (5.6,0) -- node[right]{renflement terrestre} (5.6,0.99);
\draw[<->,omThm] (6.8,2.3) -- node[right]{rayon de Fresnel} (6.8,1.78);
\node[omDef] at (3.0,3.1) {ligne de visée};
\fill[gray!50] (4.9,0.99) rectangle (5.3,1.3); \node[font=\scriptsize] at (4.4,1.2) {obstacle};
\end{tikzpicture}

\omcaption{Un saut micro-ondes, schématiquement~: les pylônes doivent élever la ligne de visée au-dessus du \omterm{def:nw:long-haul-wireless:fresnel}{renflement terrestre} au milieu du saut, en plus des obstacles et
de la première \omterm{def:nw:long-haul-wireless:fresnel}{zone de Fresnel} autour de la droite. Échelle verticale exagérée.}\label{fig:nw:long-haul-wireless:hop}
\end{omfigure}

\begin{definition}[Bilan de liaison, marge d'évanouissement]\label{def:nw:long-haul-wireless:budget}
Le \emph{bilan de liaison}\index{bilan de liaison} (link budget) additionne la puissance de l'émetteur et les gains des antennes, puis soustrait les pertes du trajet pour donner la
puissance reçue. Sa \emph{marge d'évanouissement}\index{marge d'évanouissement} (fade margin) est la puissance reçue diminuée de la puissance minimale à laquelle le récepteur décode encore le
signal au taux d'erreur requis~: la perte que la liaison peut absorber avant de défaillir.
\end{definition}

Le bilan découle de la physique une fois l'équipement fixé. En espace libre, la perte vaut $92.45 + 20\log_{10} f + 20\log_{10} d$ décibels ($f$ en
gigahertz, $d$ en kilomètres)~; une antenne parabolique de diamètre $D$ procure un gain de $10\log_{10}[\eta(\pi D f/c_0)^2]$. Avec l'équipement du modèle (30~dBm
émis, paraboles de 1.8~mètre de diamètre à 60~\% de rendement, \qty{3}{\decibel} d'autres pertes, un seuil de récepteur de $-70$~dBm, autant d'hypothèses),
un saut de 60 kilomètres conserve entre 31 et \qty{54}{\decibel} de marge (\cref{tab:nw:long-haul-wireless:budget}). Les fréquences plus élevées gagnent davantage dans
les paraboles qu'elles n'en perdent en espace libre~; la pluie reprend ce gain.

\section{Météo, évanouissement par la pluie et disponibilité}

\begin{definition}[Évanouissement par la pluie, disponibilité de la liaison]\label{def:nw:long-haul-wireless:rain}
L'\emph{évanouissement par la pluie}\index{evanouissement@évanouissement par la pluie} (rain fade) est l'atténuation d'un signal radio par les gouttes de pluie le long de son trajet, qui croît avec l'intensité de la pluie et fortement
avec la fréquence. La \emph{disponibilité de la liaison}\index{disponibilité de la liaison} (link availability) est la fraction du temps pendant laquelle elle achemine du trafic au
taux d'erreur requis, généralement exprimée par an.
\end{definition}

Recommendation ITU-R P.838-3 donne l'atténuation spécifique de la pluie sous la forme d'une loi de puissance de l'intensité de la pluie, $\gamma = k R^{\alpha}$ décibels par kilomètre
pour $R$ millimètres par heure, avec des coefficients tabulés en fonction de la fréquence et de la polarisation. À \qty{11}{\giga\hertz} (polarisation horizontale)
$k = 0.01772$ et $\alpha = 1.2140$~: une tempête de \qty{50}{\milli\metre\per\hour} coûte \qty{2.05}{\decibel} par kilomètre. À \qty{6}{\giga\hertz},
$k = 0.0007056$ et $\alpha = 1.59$~; à \qty{80}{\giga\hertz}, $k = 1.1704$ et $\alpha = 0.7115$.
La \Cref{fig:nw:long-haul-wireless:gamma} montre l'écart~: sous une pluie légère de \qty{2.5}{\milli\metre\per\hour}, un facteur d'environ 740 entre 6 et \qty{80}{\giga\hertz}.

\begin{omfigure}
\begin{tikzpicture}
\begin{semilogyaxis}[width=11cm,height=6.2cm,xmin=0,xmax=150,ymin=0.0005,ymax=60,xlabel={taux de pluie (mm/h)},
  ylabel={atténuation spécifique (dB/km)},tick label style={font=\footnotesize},label style={font=\small},grid=major,
  grid style={gray!25},legend pos=south east,legend style={font=\footnotesize},table/col sep=comma]
\addplot[omDef,thick,mark=*,mark size=1.3pt] table[x=rain,y=g6]{figdata/networks/14-long-haul-wireless/gamma.csv};
\addplot[omThm,thick,mark=square*,mark size=1.3pt] table[x=rain,y=g11]{figdata/networks/14-long-haul-wireless/gamma.csv};
\addplot[omMeth,thick,mark=triangle*,mark size=1.6pt] table[x=rain,y=g18]{figdata/networks/14-long-haul-wireless/gamma.csv};
\addplot[gray,thick,mark=diamond*,mark size=1.6pt] table[x=rain,y=g80]{figdata/networks/14-long-haul-wireless/gamma.csv};
\legend{6 GHz,11 GHz,18 GHz,80 GHz}
\end{semilogyaxis}
\end{tikzpicture}

\omcaption{Atténuation spécifique de la pluie en fonction de l'intensité de la pluie à quatre fréquences, polarisation horizontale, d'après les coefficients de la recommandation UIT-R P.838-3. Données~:
\texttt{fig\_wireless.py}~; coefficients dans \texttt{code/firm/radiolink/data/p838\_3.csv}.}\label{fig:nw:long-haul-wireless:gamma}
\end{omfigure}

\begin{proposition}[La pluie qui brise un bond]\label{prop:nw:long-haul-wireless:critical}
Un bond dont la \omterm{def:nw:long-haul-wireless:budget}{marge d'évanouissement} vaut $M$ décibels tombe en panne lorsqu'une pluie de taux $R$ couvre une longueur $L$ de celui-ci avec $k R^{\alpha} L > M$, c'est-à-dire au-delà du
taux critique
\[
  R^{*} = \left(\frac{M}{k L}\right)^{1/\alpha}.
\]
\end{proposition}

\begin{proof}
L'atténuation de la pluie sur la longueur mouillée vaut $\gamma L = k R^{\alpha} L$~; le bond tombe en panne lorsqu'elle dépasse la marge. En résolvant en $R$, on obtient
$R^{*}$.
\end{proof}

Un bond de 60~kilomètres à \qty{11}{\giga\hertz} conserve \qty{36.4}{\decibel} de marge dans le modèle~: une pluie de \qty{18.4}{\milli\metre\per\hour} le long de
toute sa longueur suffirait à le rompre, mais une cellule orageuse de dix kilomètres de large en exige \qty{80}{\milli\metre\per\hour}. À \qty{6}{\giga\hertz}, le même
bond résiste à \qty{64}{\milli\metre\per\hour} partout~; à \qty{80}{\giga\hertz}, il tombe en panne sous une bruine, ce qui explique que les longs bonds utilisent les bandes basses et que les courts
puissent utiliser les plus hautes.

\begin{table}[ht]
\centering\small
\begin{tabular}{@{}>{\centering\arraybackslash}p{1.6cm}>{\centering\arraybackslash}p{1.8cm}>{\centering\arraybackslash}p{2.8cm}>{\centering\arraybackslash}p{1.3cm}>{\centering\arraybackslash}p{2.5cm}>{\centering\arraybackslash}p{1.9cm}>{\centering\arraybackslash}p{2.0cm}@{}}
\toprule
Fréquence & Gain de l'antenne & Affaiblissement en espace libre & Reçu & \omterm{def:nw:long-haul-wireless:budget}{Marge d'évanouissement} & $R^*$, bond entier & $R^*$, cellule de 10~km \\ (GHz) & (dBi) & (dB) & (dBm) & (dB) & (mm/h) & (mm/h) \\
%
\midrule
6 & 38.9 & 143.6 & $-38.9$ & 31.1 & 63.5 & 196.0 \\
11 & 44.1 & 148.8 & $-33.6$ & 36.4 & 18.4 & 80.4 \\
18 & 48.4 & 153.1 & $-29.3$ & 40.7 & 8.1 & 42.3 \\
80 & 61.4 & 166.1 & $-16.4$ & 53.6 & 0.7 & 8.5 \\
\bottomrule
\end{tabular}
\caption{\omterm{def:nw:long-haul-wireless:budget}{Bilan de liaison} d'un bond de 60~kilomètres avec le matériel du modèle (30~dBm, deux antennes paraboliques de 1.8~mètre à 60~\% de rendement, \qty{3}{\decibel}
d'autres pertes, seuil de $-70$~dBm), et taux de pluie critique de la \cref{prop:nw:long-haul-wireless:critical} pour une pluie sur tout le bond
et pour une cellule de dix kilomètres de large. Données~: \texttt{nw\_wireless.budget\_rows()}.}\label{tab:nw:long-haul-wireless:budget}
\end{table}

\omcode{code/firm/radiolink/firm_radiolink.py}{83}{97}{Pluie sur un bond~: la loi de puissance de l'UIT-R sur la longueur mouillée, et le taux auquel elle consomme
la marge d'évanouissement.}\label{lst:nw:long-haul-wireless:rain}

\begin{omfigure}
\begin{tikzpicture}
\begin{axis}[width=11cm,height=6cm,xmin=5,xmax=80,ymin=0,ymax=440,xlabel={longueur du bond (km)},
  ylabel={taux de pluie critique (mm/h)},tick label style={font=\footnotesize},label style={font=\small},grid=major,
  grid style={gray!25},legend columns=3,legend style={font=\footnotesize,at={(0.5,-0.25)},anchor=north,draw=none},table/col sep=comma]
\addplot[omDef,thick,no markers] table[x=d_km,y=r6]{figdata/networks/14-long-haul-wireless/critical.csv};
\addplot[omThm,thick,no markers] table[x=d_km,y=r11]{figdata/networks/14-long-haul-wireless/critical.csv};
\addplot[omMeth,thick,no markers] table[x=d_km,y=r18]{figdata/networks/14-long-haul-wireless/critical.csv};
\legend{6 GHz,11 GHz,18 GHz}
\end{axis}
\end{tikzpicture}

\omcaption{Le taux de pluie auquel un bond perd sa \omterm{def:nw:long-haul-wireless:budget}{marge d'évanouissement}, en fonction de la longueur du bond, pour une cellule orageuse de dix kilomètres de large sur le bond (tout le bond
lorsqu'il mesure moins de dix kilomètres), avec le matériel du modèle. Les longs bonds perdent leur marge au profit de l'espace libre~; les fréquences plus élevées la perdent au profit de la pluie. Données~: \texttt{fig\_wireless.py},
\texttt{nw\_wireless.critical\_curve()}.}\label{fig:nw:long-haul-wireless:critical}
\end{omfigure}

La disponibilité découle de la fréquence de dépassement du taux critique, qui est une propriété du climat le long de la route, publiée dans des statistiques de taux de pluie
que ce chapitre n'utilise pas. Avec un modèle stylisé dans lequel il pleut 5~\% du temps avec un taux moyen de
\qty{4}{\milli\metre\per\hour} (distribué exponentiellement), un bond de 59~kilomètres à \qty{11}{\giga\hertz} arrosé de bout en bout est hors service 0.047~\% de
l'année, soit environ quatre heures~; vingt bonds de ce type tombant en panne indépendamment mettraient la route sur la fibre pendant environ 0.9~\% du temps. Le modèle est grossier~;
sa structure ne l'est pas~: la disponibilité se multiplie le long d'une chaîne, et une route n'est jamais plus sèche que son bond le plus arrosé.

\section{Ondes millimétriques et lasers}

\begin{definition}[Liaison en ondes millimétriques, liaison optique en espace libre]\label{def:nw:long-haul-wireless:mmw}
Une \emph{liaison en ondes millimétriques}\index{liaison en ondes millimétriques} (millimetre-wave link) est une liaison radio à des fréquences d'environ 30 à 300~GHz (longueurs d'onde de 10 à 1
millimètre), à larges canaux et à bonds courts. Une \emph{liaison optique en espace libre}\index{liaison optique en espace libre} (free-space optical link, FSO) transporte les données sur un faisceau laser
à travers l'air entre deux terminaux~; le brouillard et les fortes pluies le bloquent, et elle exige une visibilité directe, claire et stable.
\end{definition}

À l'intérieur d'une zone métropolitaine, le compromis s'inverse~: les bonds sont courts, la pluie a peu de trajet sur lequel agir, et la bande passante compte. Les routes de l'ICE dans le New Jersey
entre Mahwah et Secaucus, exploitées avec Anova Financial Networks, combinent ondes millimétriques et optique en espace libre sur des canaux radio de 5~gigabits~;
entre Mahwah et Carteret, elles utilisent plusieurs canaux de 1~gigabit. Les deux technologies tombent en panne différemment (la pluie pour les ondes millimétriques,
le brouillard pour les lasers), si bien qu'une route qui porte les deux ne perd les deux que lorsque la pluie et le brouillard surviennent ensemble.

\section{Ondes courtes à travers les océans}

\begin{definition}[Propagation par onde ionosphérique]\label{def:nw:long-haul-wireless:sky}
La \emph{propagation par onde ionosphérique}\index{propagation par onde ionosphérique} (skywave) transporte des ondes radio de quelques à quelques dizaines de mégahertz au-delà de l'horizon par réflexion sur
l'ionosphère, une ou plusieurs fois, de sorte qu'une seule liaison puisse franchir un océan.
\end{definition}

La géodésique de Secaucus NY4 à Slough LD4 mesure \qty{5551}{\kilo\metre}~: \qty{18.52}{\milli\second} de lumière dans le vide, contre
\qty{27.07}{\milli\second} dans du verre en ligne droite et \qty{29.48}{\milli\second} pour la route sous-marine publiée la plus rapide du chapitre~13. Une seule réflexion
sur une couche haute de 300~kilomètres allonge un trajet de 5\,551~kilomètres de seulement environ 0.6~\% dans l'approximation de la Terre plate, de sorte qu'une liaison par onde ionosphérique peut
battre n'importe quel câble d'environ dix millisecondes, hors matériel, avec une bande passante bien inférieure à celle d'un canal hertzien et une disponibilité qui suit l'ionosphère, heure par heure.

\begin{dated}{2026-09}{Ondes courtes pour le trading}\label{dat:nw:long-haul-wireless:hf}
Les médias spécialisés ont signalé en juillet 2023 que des traders à haute fréquence utilisaient depuis plusieurs années des liaisons par onde ionosphérique pour transmettre des données de trading entre des bourses américaines et
étrangères, dans le cadre d'autorisations expérimentales, et qu'une Shortwave Modernization Coalition souhaitait ouvrir la bande \qtyrange{2}{25}{\mega\hertz}
à un usage fixe, longue distance et non vocal. La FCC a demandé des observations sur la requête de la coalition (RM-11953) par un avis public du 30 juin 2023.
\end{dated}

\section{Licences, pylônes, bande passante et contenu transmis}

Une route hertzienne est publique d'une façon qu'une route en fibre ne l'est pas~: aux États-Unis, ses trajets sont soumis à licence, et Laughlin, Aguirre et Grundfest ont
utilisé les registres de licences de la FCC pour recenser au moins quinze réseaux Chicago--New York, leurs pylônes et leurs longueurs. Le spectre est étroit~:
les réseaux privés y utilisent surtout des canaux de 30 et 40~mégahertz dans les bandes de 6 et 11~gigahertz, à environ 140 à 190~mégabits par seconde par
canal selon leur propre décompte. Les services vendus par les plateformes sont plus étroits encore~: de 1 à 10~mégabits de bande passante privée sur le service
de l'ICE à Toronto, 10~mégabits sur le réseau Zurich--Francfort de SIX.

Ce qui est transmis est donc choisi~: les quelques champs des quelques instruments qui font bouger un autre marché (le meilleur prix d'un contrat à terme sur indice, une
transaction), dans des formats compacts, le reste passant par la fibre. La latence propre de la radio compte~: Laughlin et ses coauteurs décrivent des équipements
du marché en 2010 qui ajoutaient de quelques dizaines à plusieurs milliers de microsecondes par saut, surtout pour la \omterm{def:nw:long-haul-fibre:fec}{correction d'erreurs directe} (FEC) et l'entrelacement, et estiment de quelques à une vingtaine
de microsecondes par saut pour les réseaux alors en cours de licence. Une route de vingt sauts multiplie chaque microseconde par vingt.

\section{Un orage à travers la route}

Le modèle fait passer un orage sur la route Aurora--Carteret~: vingt sauts de \qty{59}{\kilo\metre} sur le grand cercle (\qty{1180.9}{\kilo\metre}), une
cellule pluvieuse large de \qty{30}{\kilo\metre} traversant la ligne à \qty{48}{\kilo\metre\per\hour}, dont l'intensité décroît linéairement depuis un maximum à
son centre jusqu'à zéro sur ses bords. Tant que chaque saut conserve sa marge, la route fonctionne à la latence publiée de \qty{3.982}{\milli\second}~; lorsqu'un saut tombe, elle
bascule sur la fibre, que le modèle traite comme un trajet plus long de 20~\% que la géodésique, soit \qty{6.911}{\milli\second}. Tout, hormis la géométrie,
les coefficients et la latence publiée de la route hertzienne, relève d'une hypothèse.

\begin{omfigure}
\begin{tikzpicture}
\begin{axis}[width=11cm,height=5.8cm,xmin=0,xmax=37,ymin=3.5,ymax=7.5,xlabel={minutes depuis l'arrivée de la cellule sur la route},
  ylabel={latence de la route (ms, sens unique)},tick label style={font=\footnotesize},label style={font=\small},grid=major,
  grid style={gray!25},legend columns=3,legend style={font=\footnotesize,at={(0.5,-0.25)},anchor=north,draw=none},table/col sep=comma]
\addplot[omDef,thick,const plot] table[x=minute,y=p50]{figdata/networks/14-long-haul-wireless/storm.csv};
\addplot[omThm,thick,dashed,const plot] table[x=minute,y=p100]{figdata/networks/14-long-haul-wireless/storm.csv};
\addplot[omMeth,thick,dotted,const plot] table[x=minute,y=p150]{figdata/networks/14-long-haul-wireless/storm.csv};
\legend{pic 50 mm/h,pic 100 mm/h,pic 150 mm/h}
\end{axis}
\end{tikzpicture}

\omcaption{Simulation, et non une mesure~: la latence unidirectionnelle de la route Aurora--Carteret, minute par minute, pendant qu'une cellule pluvieuse large de \qty{30}{\kilo\metre}
traverse l'un de ses sauts à \qty{11}{\giga\hertz} à \qty{48}{\kilo\metre\per\hour}, pour trois intensités de pointe~: la latence publiée de la route hertzienne,
\qty{3.982}{\milli\second}, ou la fibre à une valeur supposée de \qty{6.911}{\milli\second} tant qu'un saut est coupé. Données~: \texttt{fig\_wireless.py},
\texttt{nw\_wireless.storm()}.}\label{fig:nw:long-haul-wireless:storm}
\end{omfigure}

Avec un pic de \qty{100}{\milli\metre\per\hour}, la route passe 21 des 38 minutes de l'orage sur la fibre, soit \qty{2.93}{\milli\second} de plus~; avec
\qty{50}{\milli\metre\per\hour}, 5 minutes~; avec \qty{150}{\milli\metre\per\hour}, 26. Les mêmes orages à \qty{6}{\giga\hertz} coûtent 0, 0 et 7
minutes~; à \qty{18}{\giga\hertz}, 22, 29 et 31. Pour les firmes qui se livrent à cette course, ces minutes sont les après-midi de l'étude de Shkilko et Sokolov~; pour
leurs contreparties, ce sont les moments où le marché est plus clément.

\omcode{code/firm/radiolink/firm_radiolink.py}{130}{150}{L'orage~: la corde de la cellule à travers la route, la longueur mouillée et l'intensité de pluie sur chaque saut
sous elle, et la latence de la route, radio ou fibre, minute par minute.}\label{lst:nw:long-haul-wireless:storm}

\begin{method}[Évaluer une route hertzienne]\label{met:nw:long-haul-wireless:eval}
\begin{enumerate}
\item À partir des registres de licences ou auprès du fournisseur, obtenir les pylônes~: nombre de sauts, longueurs, fréquences. Calculer le \omterm{def:nw:the-north-american-map:factor}{facteur de tracé} du
  trajet (chapitre~10).
\item Pour chaque saut, calculer la géométrie et, avec les caractéristiques d'équipement du fournisseur, la \omterm{def:nw:long-haul-wireless:budget}{marge d'évanouissement}~; repérer les sauts dont l'intensité
  de pluie critique est la plus faible.
\item Combiner avec le climat le long de la route pour estimer la disponibilité~; demander au fournisseur sa disponibilité mesurée et son historique de pannes.
\item Chiffrer la fibre de secours et concevoir la bascule~: qui détecte l'évanouissement, à quelle vitesse le trafic se déplace, et ce que font les stratégies pendant ce temps.
\item Décider de ce que transporte la radio~: les quelques messages dont la vitesse justifie le coût.
\end{enumerate}
\end{method}

\section{Tutoriel~: un saut, la pluie et un orage}\label{tut:nw:long-haul-wireless:tut}

\begin{tutorial}
\textbf{Objectif.} Calculer la géométrie et le bilan d'un saut, ses limites de pluie par fréquence, et le coût d'un orage en minutes sur la fibre.
\textbf{État final~:} \cref{fig:nw:long-haul-wireless:gamma,fig:nw:long-haul-wireless:critical,fig:nw:long-haul-wireless:storm} et
\cref{tab:nw:long-haul-wireless:geometry,tab:nw:long-haul-wireless:budget}.
\begin{tutsteps}
\item \textbf{Géométrie.} \texttt{firm\_radiolink.bulge\_m}, \texttt{fresnel\_m} et \texttt{tower\_m} (\cref{prop:nw:long-haul-wireless:geometry})~;
  \texttt{nw\_wireless.geometry\_rows()} remplit le \cref{tab:nw:long-haul-wireless:geometry}.
\item \textbf{Bilan.} \texttt{Hop}, \texttt{rx\_dbm} et \texttt{fade\_margin\_db}~; \texttt{budget\_rows()} remplit
  le \cref{tab:nw:long-haul-wireless:budget}.
\item \textbf{Pluie.} Les coefficients de la recommandation P.838-3 sont des données (\texttt{data/p838\_3.csv}), chaque ligne accompagnée de sa référence au registre~;
  \texttt{critical\_rain} (\cref{lst:nw:long-haul-wireless:rain}) met en œuvre la \cref{prop:nw:long-haul-wireless:critical}.
\item \textbf{La route et l'orage.} \texttt{chain} place vingt sauts sur le grand cercle~; \texttt{storm\_timeline}
  (\cref{lst:nw:long-haul-wireless:storm}) fait passer la cellule sur l'un d'eux.
\end{tutsteps}
\textbf{Pour aller plus loin.} Utiliser la polarisation verticale (atténuation plus faible à ces fréquences)~; raccourcir les sauts à 40~kilomètres et ajouter
  des pylônes~; déplacer l'orage le long de la route au lieu de la traverser.
\end{tutorial}

\section{Construction~: le modèle de liaison radio}\label{bld:nw:long-haul-wireless:radiolink}

\begin{build}
\bfield{Rôle} Sauts hertziens et chaînes avec leur géométrie, leur bilan et leur pluie, et la bascule (failover) vers la fibre~: la partie aérienne des routes que les
  chapitres~15, 19 et~29 chiffrent et planifient.
\bfield{Interface} \texttt{firm\_radiolink}~: \texttt{bulge\_m}, \texttt{fresnel\_m}, \texttt{tower\_m}, \texttt{fspl\_db}, \texttt{dish\_gain\_dbi},
\texttt{Hop}, \texttt{rx\_dbm}, \texttt{fade\_margin\_db}, \texttt{coefficients}, \texttt{gamma\_db\_km}, \texttt{rain\_db}, \texttt{critical\_rain},
\texttt{chain}, \texttt{hop\_lengths}, \texttt{Storm}, \texttt{storm\_timeline}, \texttt{availability}.
\bfield{Règles} Les coefficients de pluie proviennent du fichier de données avec leur source, jamais saisis dans le code~; les caractéristiques d'équipement sont des paramètres~; un orage est une
  simulation et le dit.
\bfield{Tests d'acceptation} \texttt{code/firm/radiolink/tests/}~: les règles du milieu de saut, la perte en espace libre, un bilan et la pluie critique calculées à la main, les
  extrémités de la chaîne sur le grand cercle, un orage sec et un orage pluvieux, et la disponibilité aux extrêmes.
\bfield{Pour aller plus loin} Profils de terrain le long de chaque saut~; la méthode de l'UIT-R pour la longueur effective de trajet sous la pluie~; pluie corrélée entre
  sauts voisins.
\end{build}

\begin{omsources}
\item A.~Shkilko et K.~Sokolov, «~Every cloud has a silver lining: fast trading, microwave connectivity, and trading costs~», \emph{Journal of
Finance} 75(6) (2020).
\item Recommandation UIT-R P.838-3 (2005), modèle d'atténuation spécifique due à la pluie.
\item G.~Laughlin, A.~Aguirre et J.~Grundfest (2014), section sur les réseaux hertziens~; routes hertziennes de l'ICE, région métropolitaine du New Jersey~; FCC Order DA 23-689
(2023) et Radio World (juillet 2023) sur les ondes courtes.
\end{omsources}

\section{Exercices}

\begin{exercise}[$\star$]\label{exo:nw:long-haul-wireless:1}
Calculer la hauteur du \omterm{def:nw:long-haul-wireless:fresnel}{renflement terrestre} et le rayon de la première \omterm{def:nw:long-haul-wireless:fresnel}{zone de Fresnel} au milieu d'un saut de 40~kilomètres à \qty{11}{\giga\hertz}, avec $K = 4/3$.
\end{exercise}

\begin{exercise}[$\star$]\label{exo:nw:long-haul-wireless:2}
Calculer l'atténuation spécifique de la pluie à \qty{25}{\milli\metre\per\hour}, à 6 et à \qty{18}{\giga\hertz}.
\end{exercise}

\begin{exercise}[$\star$]\label{exo:nw:long-haul-wireless:3}
Pourquoi un réseau hertzien peut-il avoir un \omterm{def:nw:the-north-american-map:factor}{facteur de tracé} voisin de 1.01 alors qu'un tracé en fibre peine à descendre sous 1.2~?
\end{exercise}

\begin{exercise}[$\star\star$]\label{exo:nw:long-haul-wireless:4}
Doubler la longueur d'un saut ajoute combien de décibels de perte en espace libre, et que fait-il du taux de pluie critique si la pluie couvre tout le saut~?
\end{exercise}

\begin{exercise}[$\star\star$]\label{exo:nw:long-haul-wireless:5}
Pourquoi les routes métropolitaines associent-elles les ondes millimétriques aux lasers~?
\end{exercise}

\begin{exercise}[$\star\star$]\label{exo:nw:long-haul-wireless:6}
Une liaison en ondes courtes traverse l'Atlantique dix millisecondes plus vite que n'importe quel câble. Que peut-elle transporter, et que ne peut-elle pas transporter~?
\end{exercise}

\begin{exercise}[$\star\star\star$]\label{exo:nw:long-haul-wireless:7}
\emph{Programmation.} Avec \texttt{firm.radiolink}, trouver le taux de pluie maximal à partir duquel la tempête du modèle met pour la première fois la route à \qty{11}{\giga\hertz} sur fibre.
\end{exercise}

\begin{exercise}[$\star\star\star$]\label{exo:nw:long-haul-wireless:8}
\emph{Trouver la faille.} «~Notre route compte vingt sauts, chacun disponible 99.95~\% de l'année, donc la route est disponible 99.95~\% de l'année.~»
\end{exercise}

\section{Problème~: La tempête entre Chicago et le New Jersey}

\begin{problem}[{Problème du week-end --- ce que coûte une tempête à une route radio}]\label{pb:nw:long-haul-wireless:1}
Une route radio de l'entreprise, d'Aurora à Carteret, compte vingt sauts d'environ \qty{59}{\kilo\metre} à \qty{11}{\giga\hertz}, fonctionne à \qty{3.982}{\milli\second}
dans un sens, et bascule sur la fibre à \qty{6.911}{\milli\second}. Utiliser l'équipement et la tempête du modèle.

\medskip\noindent\textbf{Partie I --- Un saut.}
\begin{enumerate}
\item Quels sont le \omterm{def:nw:long-haul-wireless:fresnel}{renflement terrestre} et le rayon de Fresnel au milieu du saut, ainsi que la hauteur du pylône pour un obstacle de 10~mètres~?
\item Quelles sont la perte en espace libre du saut, la puissance reçue et la \omterm{def:nw:long-haul-wireless:budget}{marge d'évanouissement}~?
\item À quel taux de pluie le saut perd-il sa marge si la pluie couvre tout le saut~?
\item Et si la pluie ne couvre que dix kilomètres de celui-ci~?
\end{enumerate}

\medskip\noindent\textbf{Partie II --- La tempête.}
\begin{enumerate}[resume]
\item Combien de temps une cellule de 30~kilomètres de large met-elle pour traverser la route à 48~kilomètres à l'heure~?
\item Avec un pic de \qty{100}{\milli\metre\per\hour}, combien de minutes la route passe-t-elle sur fibre~?
\item De combien est-elle plus lente pendant ces minutes~?
\item Comment la réponse change-t-elle à 6 et à \qty{18}{\giga\hertz}~?
\end{enumerate}

\medskip\noindent\textbf{Partie III --- L'année.}
\begin{enumerate}[resume]
\item Avec le modèle de dépassement étiqueté, combien de temps un saut est-il hors service en une année si la pluie le couvre de bout en bout~?
\item Que donne cela pour vingt sauts indépendants~?
\item Pourquoi l'hypothèse d'indépendance est-elle douteuse~?
\item Que donneraient des sauts plus courts~?
\end{enumerate}

\medskip\noindent\textbf{Partie IV --- Le verdict.}
\begin{enumerate}[resume]
\item Énoncer le \emph{résultat nommé}~: le taux de pluie à partir duquel un saut de 59~kilomètres à \qty{11}{\giga\hertz} perd sa \omterm{def:nw:long-haul-wireless:budget}{marge d'évanouissement}, et les minutes de
  latence fibre seule que coûte une tempête à \qty{100}{\milli\metre\per\hour} traversant la route.
\item Qu'ont constaté Shkilko et Sokolov quant à l'évolution des coûts de transaction pendant de tels épisodes~?
\item Pourquoi une entreprise conserverait-elle une route à 6~gigahertz et une route à 18~gigahertz~?
\item Que doivent faire les systèmes de l'entreprise lorsque la route bascule sur la fibre~?
\item Laquelle des entrées du modèle remplacerait-on en premier par des données réelles, et d'où les tirerait-on~?
\item Que change la tempête pour une entreprise qui ne fait pas du tout de course à la latence~?
\item La latence de la sauvegarde sur fibre vaut-elle la peine d'être réduite~? Que faudrait-il faire (chapitre~13)~?
\item En une phrase~: que fait la météo à une course à la latence~?
\end{enumerate}
\end{problem}

\section{Questions d'entretien}

\begin{interviewq}[$\star$ \iqroles{développeur}]\label{iq:nw:long-haul-wireless:1}
Pourquoi les micro-ondes sont-elles plus rapides que la fibre entre deux villes, et que coûtent-elles en contrepartie~?
\end{interviewq}

\begin{interviewq}[$\star\star$ \iqroles{développeur}]\label{iq:nw:long-haul-wireless:2}
Qu'est-ce qui limite la longueur d'un saut hertzien~?
\end{interviewq}

\begin{interviewq}[$\star\star$ \iqroles{développeur, trader}]\label{iq:nw:long-haul-wireless:3}
Il pleut en Pennsylvanie. Que devient votre stratégie Chicago--New Jersey, et que doivent faire vos systèmes~?
\end{interviewq}

\begin{interviewq}[$\star\star$ \iqroles{trader, chercheur}]\label{iq:nw:long-haul-wireless:4}
Comment utiliseriez-vous les interruptions de réseaux hertziens causées par la météo pour étudier l'effet de la vitesse sur la qualité de marché~?
\end{interviewq}

\begin{interviewq}[$\star\star$ \iqroles{développeur}]\label{iq:nw:long-haul-wireless:5}
Que transmettriez-vous sur une liaison hertzienne de 10~mégabits entre deux bourses, et que laisseriez-vous sur la fibre~?
\end{interviewq}

\begin{interviewq}[$\star\star\star$ \iqroles{développeur, chercheur}]\label{iq:nw:long-haul-wireless:6}
Comparez les micro-ondes, les ondes millimétriques, les lasers et les ondes courtes pour relier deux sites de négociation distants de 50, 1\,000 et 6\,000 kilomètres.
\end{interviewq}
